\section{Discussion}
\label{sec:discussion}

\subsection{Answering the research question}

Cloud-native security telemetry does detect all five attack classes, quickly and
at a false-positive rate a small team can absorb --- but the qualifier matters
more than the headline. Three of the five attacks were first identified by logs
the \emph{ground segment} produces, not by anything a cloud provider generates on
its own. The cloud tier supplied the pipeline, the correlation, the findings
store and the response automation; it did not supply the evidence. An
organisation that adopts cloud-native security tooling and instruments only what
that tooling emits by default will detect credential abuse and data
exfiltration, and will be blind to the attacks that reach the vehicle.

The honest form of the answer is therefore conditional: cloud-native security
telemetry is effective for ground segments \emph{if the ground segment is
instrumented as a first-class log source}. That conclusion is unglamorous and, in
our view, the most useful thing this study produces.

\subsection{Detection is not a substitute for prevention}

The ablation makes this concrete in a way that a narrative argument cannot.
Removing frame authentication does not reduce the detection rate at all --- every
forged command is still detected, promptly --- yet it raises the number of
critical commands executing on the spacecraft from \AblFullCrit{} to
\AblNoSDLSCrit{}. A perfect detector and a \SI{4}{\second} pipeline do not undo a
thruster firing.

This is a general property of systems with physical actuators, and it argues for
a specific division of labour. Preventive controls belong as close to the
actuator as possible --- on the vehicle, where the anti-replay window lives, not
in the analytics tier. Detection exists to catch what prevention does not cover
and to bound the duration of an intrusion, not to replace it. In a ground segment
the distinction is unusually sharp, because the window in which prevention can
act is the pass, and the pass is over in seven minutes.

\subsection{Cryptographic integrity and content analytics are complementary}

Attack A3 was designed to separate two defences that are often treated as
alternatives. Its first phase injects implausible values with a forged tag: the
cryptographic check catches it immediately. Its second phase first disables frame
authentication and then freezes a degrading measurement at a plausible constant.
After that point the tag check has nothing to say --- every frame is
``correctly'' unauthenticated --- and the only detection left is the one that
reasons about whether the measurement makes physical sense.

The measurement adds a wrinkle we did not anticipate. Once authentication is
off, the physics rule also inherits the link's own noise, because the ground
segment can no longer separate a corrupted frame from an injected one. The
fallback detection is therefore least precise exactly when it becomes the only
one --- a property worth knowing before treating content analytics as a
substitute for provenance.

Neither defence subsumes the other. Cryptography establishes provenance and says
nothing about content; physics analytics evaluates content and says nothing about
provenance. A ground segment that has only one of them has a documented blind
spot, and A3 traverses it in two moves.

\subsection{Where the false positives came from}

Every false positive in the campaign came from a thresholded or behavioural rule,
and each traced to a benign anomaly that a real ground segment produces every
week. This has a practical implication for anyone building such a pack: the
policy rules --- refused commands, refused frames, stopped audit trails --- are
close to free, and should be written first. They are precise because the events
they match genuinely do not occur in nominal operations, and that property comes
from the ground segment's small, well-defined operational envelope rather than
from clever rule authoring.

The behavioural rules earn their noise, though. R02, the unprofiled-origin rule,
was the first detector for two of the five attacks. In a general enterprise,
origin profiling is a notoriously noisy control; in a ground segment, where the
operator population is a handful of named people working from a handful of known
networks, it is one of the strongest signals available. Threat detection content
does not transfer between domains unchanged, and this is a case where the domain
makes a weak enterprise rule into a strong one.

\subsection{Containment degrades the observability of later attacks}

The A2 latency distribution in Section~\ref{sec:results} is bimodal for a reason
that has nothing to do with the adversary. Automated containment blocks a
principal, and in this architecture a block is permanent for the remainder of the
campaign. When a false positive blocks a benign operator --- which happens at the
false-positive rate the pack already reports --- that account stops being usable,
and the mission-control API refuses any later session for it \emph{silently},
without emitting an authentication event. An adversary who subsequently targets
that account therefore produces no logon for the behavioural rules to evaluate.
Detection does not fail, but it falls back to a slower path: in our campaign,
from \ResTTDMedianAtwo{}\,s on the session to as much as \ResTTDMaxAtwo{}\,s on
the radio-frequency phase, which can only be observed during a pass.

This is worth stating plainly because it inverts an assumption that is easy to
make. Automated response is normally evaluated on what it prevents, and
Section~\ref{sec:results} reports that in those terms. But a response action also
changes what the telemetry contains, and a control that removes an actor from the
system removes that actor's events along with it. The effect is not confined to
this testbed: any architecture that disables accounts automatically and logs
refusals at a different fidelity than successes will have the same blind spot.
Two mitigations follow directly --- expire automated blocks rather than leaving
them indefinite, and emit an explicit event when an action is refused because the
principal is blocked, so that the refusal is itself a detectable signal. We did
not implement either here; measuring them is future work.

\subsection{Implications for practitioners}

\begin{enumerate}
\item \textbf{Ship ground-segment logs.} Authentication, telecommand submission,
uplink verdicts and frame integrity are the evidence for three of five attacks.
Nothing generates them for you.
\item \textbf{Put the anti-replay window on the vehicle and verify it.} It is the
only control in the study that both prevents an attack and produces its alert.
\item \textbf{Choose your query cadence deliberately.} Moving batch-tier latency
from \SI{900}{\second} to \SI{30}{\second} improved campaign mean latency by more
than any rule change we made, at no cost in false positives.
\item \textbf{Exclude known-good service principals by identity, never by
threshold.} The weekly reprocessing job and the exfiltration attack are
indistinguishable on volume; raising the threshold past the job also raises it
past the attack.
\item \textbf{Automate containment, not remediation.} Containment cut impact by
roughly four fifths in our data. Automating anything that commands the vehicle
would make a false positive into a mission outage.
\item \textbf{Budget for the platform, not the volume.} At single-mission scale
the cost is a floor, not a function of telemetry. Pool the analytics tier across
missions rather than reducing what you collect.
\end{enumerate}
